\begin{abstract}
This dissertation will investigate methods for identifying copyrighted audio within short-form and transformed media content, a problem that is becoming increasingly relevant across platforms such as TikTok, YouTube Shorts, and Instagram Reels. Existing audio fingerprinting systems frequently fail to recognise excerpts that have been pitch-shifted, time-stretched, truncated, or mixed, limiting their suitability for modern user-generated content. The proposed system, \textbf{WaveID}, will explore the use of deep learning–based audio embeddings and contrastive learning techniques to achieve robust audio traceability under such transformations. The research will utilise publicly available and ethically cleared datasets, including the Free Music Archive and GTZAN, supplemented with synthetic audio augmentations designed to replicate real-world editing behaviours. WaveID will be evaluated using quantitative metrics, robustness testing, and scalability analysis to determine its effectiveness under varying distortion conditions. The intended contribution is a reproducible framework for resilient audio identification that may support future work in digital rights management, automated content verification, and music information retrieval.
\end{abstract}
